% =============================================================================
% apendiceA.tex — PROGRAMA DE CALCULO
% =============================================================================

\chapter{PROGRAMA DE CÁLCULO}\label{ap:codigo-fonte}

O memorial apresentado do \autoref{cap:predimensionamento} ao
\autoref{cap:detalhamento} foi integralmente produzido por um programa
desenvolvido para este trabalho. Este apêndice descreve a sua arquitetura, o
fluxo de execução, o ambiente de reprodução e a forma de acesso ao código.

Optou-se por não reproduzir aqui a listagem integral. O pacote reúne \nLinhasCodigo{}
linhas distribuídas por \nArquivosCodigo{} arquivos, volume que ocuparia
várias vezes a extensão do texto e cuja leitura sequencial em papel não é praticável. O
código está publicamente disponível, conforme a
\autoref{sec:ap-disponibilidade}, em forma navegável e com o histórico completo
de alterações.

\section{ARQUITETURA}\label{sec:ap-arquitetura}

O programa está organizado em pacotes de responsabilidade única. Os dados de
entrada e as grandezas derivadas concentram-se em um pacote; cada família de
verificações ocupa o seu; e a geração de saídas é isolada do cálculo. Nenhum
módulo de verificação escreve arquivos, e nenhum módulo de saída executa
cálculo. A \autoref{tab:ap-modulos} descreve a divisão; ela é gerada na
execução, pela contagem dos próprios arquivos, e a suíte de testes aparece em
linha separada, fora do total do programa. A contagem de linhas
inclui comentários e documentação embutida, que representam parcela expressiva
do total, por registrarem a proveniência normativa de cada valor.

% tab_ap_modulos.tex — GERADA por mlc_projeto/output/apendice_programa.py (issue 121). NÃO EDITAR À MÃO.
\tabela{Pacotes do programa e respectivas responsabilidades}{\footnotesize
\begin{tabular}{lcrp{5.6cm}}
\toprule
    \textbf{Pacote} & \textbf{Arq.} & \textbf{Linhas} & \textbf{Responsabilidade} \\
\midrule
    \texttt{core} & 7 & 3.519 & Geometria paramétrica, material, ações, vento, proveniência normativa e validação de domínio \\
    \texttt{predim} & 4 & 913 & Varredura paramétrica do espaço de geometrias \\
    \texttt{checks/elu} & 30 & 11.079 & Verificações de estado-limite último de elementos e sistemas \\
    \texttt{checks/els} & 4 & 1.005 & Verificações de estado-limite de serviço \\
    \texttt{connections} & 15 & 5.657 & Ligações, chapas metálicas, bases e ancoragem \\
    \texttt{output} & 18 & 7.927 & Memorial, tabelas, figuras, desenhos, apêndices gerados, validação e custo \\
    \texttt{opensees} & 4 & 962 & Modelos de elementos finitos independentes (viga, diafragma) \\
    \texttt{visualization} & 1 & 518 & Visualização tridimensional do modelo \\
    \texttt{main.py} & 1 & 2.938 & Orquestração do \textit{pipeline} e dados de projeto \\
    \texttt{\_\_init\_\_} & 3 & 102 & Inicialização dos pacotes \\
\midrule
    \textbf{Total do programa} & \textbf{87} & \textbf{34.620} & \\
\addlinespace
    \texttt{tests} & 112 & 10.765 & Suíte de testes, fora do total do programa \\
\bottomrule
\end{tabular}}{Elaborado pelo autor (2026)}{ap-modulos}{}{}


Três módulos do pacote \texttt{core} merecem menção, por sustentarem decisões
discutidas no corpo do trabalho. O módulo de geometria calcula altura, área,
inércia e módulo resistente em qualquer posição do vão, de modo que uma
alteração de seção se propague a todo o carregamento e a todas as verificações.
O módulo de proveniência normativa registra, para cada parâmetro, a cláusula de
origem e a natureza, e é dele que decorre a declaração dos \nNDPExt{} parâmetros
de determinação nacional da \autoref{tab:ndp}. O módulo de validação impõe
pré-condições de domínio anteriores a qualquer verificação, impedindo que uma
geometria sem sentido físico produza taxa de utilização e veredito de aprovação.

\section{FLUXO DE EXECUÇÃO}\label{sec:ap-fluxo}

A execução percorre fases sucessivas, cada uma consumindo os resultados da
anterior, conforme a \autoref{tab:ap-fases}. As fases 1 e 6 são opcionais e não
integram a execução padrão, e a fase 4 inclui a auditoria de coerência do modelo
de apoio horizontal, cujo resultado consta da \autoref{tab:rigidez-apoio}.

\tabela
{Fases da execução do programa}
{\footnotesize
\begin{tabular}{lp{8.4cm}}
\toprule
    \textbf{Fase} & \textbf{Conteúdo} \\
\midrule
    0 & Dados de entrada: geometria, material, durabilidade, peso próprio por parcelas, ações de vento e combinações \\
    1 & Varredura paramétrica, quando solicitada \\
    2 & Estados-limites últimos: viga, pilares, terças, contraventamentos, nó de topo, estabilidade global, fogo, impacto e içamento \\
    3 & Estados-limites de serviço \\
    4 & Ligações, ruptura de bloco, fendilhamento, chapas metálicas, sucção, pilar de frontão, ancoragem e detalhe das bases \\
    5 & Memorial em LaTeX, planilha, figuras, desenhos cotados e apêndices B e C desta monografia \\
    6 & Visualização tridimensional, quando solicitada \\
\bottomrule
\end{tabular}}
{Elaborado pelo autor (2026)}
{ap-fases}
{}
{}

A varredura global apresentada na \autoref{sec:varredura-global} não constitui
fase do \textit{pipeline}, e sim um invólucro que o executa repetidamente. Para cada
ponto da grade, substitui a geometria e chama a mesma rotina que produz o
memorial, o que dá equivalência de execução e de resultados entre a avaliação
das candidatas e a verificação da solução adotada; não se trata de validação
física, que é a da \autoref{sec:ap-verificacao}.

\section{VERIFICAÇÃO DO PROGRAMA}\label{sec:ap-verificacao}

O mesmo programa calcula, tabela e escreve as equações do memorial. Um erro de
modelagem ou de implementação se propagaria, por isso, de forma coerente a todas
as saídas, sem contradição interna que o denunciasse: é o risco de modo comum.
Testes que apenas fixam o resultado atual do próprio programa protegem contra
alterações involuntárias, mas não contra esse risco. A resposta a ele são
verificações cujo valor esperado vem de fora do programa.

A suíte reúne \nTestes{} testes automatizados, de duas classes. Os
\nTestesUnitarios{} testes unitários exercitam funções isoladas com dados
fixos; muitos comparam o resultado com o valor calculado à mão no próprio teste. Os
\nTestesIntegracao{} testes de integração executam o programa completo e
conferem relações entre as saídas: que cada dimensão de cálculo seja a mesma da
cota de desenho e da lista de peças, que cada equação do Apêndice B reproduza,
com os números impressos, o resultado impresso, e que as verificações respeitem
os critérios declarados.

A esses somam-se \nCasosReferencia{} casos de referência independentes,
reunidos na \autoref{tab:referencias}. Em cada um, o valor esperado é obtido por
caminho alheio ao programa: solução analítica fechada da resistência dos
materiais; fórmula normativa recalculada à mão, com números fixos, sem chamar a
função do programa; equilíbrio global; refino da discretização; ou modelo de
elementos finitos independente. A tabela é gerada na execução, e a execução é
interrompida se algum caso sair da tolerância.

% tab_referencias.tex — GERADA por mlc_projeto/output/validacao.py (issue 121). NÃO EDITAR À MÃO.
\tabela{Casos de referência independentes}{\footnotesize
\begin{tabular}{p{2.2cm}p{3.1cm}p{2.2cm}rrr}
\toprule
    \textbf{Classe} & \textbf{Grandeza} & \textbf{Referência} & \textbf{Esperado} & \textbf{Programa} & \textbf{Erro} \\
\midrule
    Analítico & Viga biapoiada prismática: M no meio do vão (kN·m) & $qL^2/8$ & 500 & 500 & $<10^{-6}$ \\
    Analítico & Viga biapoiada prismática: reação (kN) & $qL/2$ & 100 & 100 & $<10^{-6}$ \\
    Analítico & Viga com balanço: momento no apoio (kN·m) & $-q\,a^2/2$ & $-$20 & $-$20 & $<10^{-6}$ \\
    Analítico & Viga com balanço: reação do apoio externo (kN) & estática & 121 & 121 & $<10^{-6}$ \\
    Analítico & Carga assimétrica nas águas: reação A (kN) & estática, $q_e = 12$, $q_d = 6$ & 105 & 105 & $<10^{-6}$ \\
    Normativo & Fator $k_c$ para $\lambda_{rel} = 1{,}2$, $\beta_c = 0{,}1$ & EN 1995-1-1, (6.27) a (6.29) & 0,6005 & 0,6005 & $<10^{-6}$ \\
    Normativo & Número efetivo, 4 pinos, $a_1 = 100$, $d = 20$ & EN 1995-1-1, (8.34) & 2,742 & 2,742 & $<10^{-6}$ \\
    Normativo & Johansen, chapa central, $d = 24$, $t_1 = 148$ (kN) & EN 1995-1-1, (8.11), (8.30) e (8.32) & 37,65 & 37,65 & $<10^{-6}$ \\
    Normativo & Cone de arrancamento, $h_{ef} = 400$ mm, C25/30 (kN) & EN 1992-4, (7.2) & 356 & 356 & $<10^{-6}$ \\
    Equilíbrio & Maior desequilíbrio $|R_A + R_B - \sum q|$, 29 combinações (kN) & estática & 0 & 0 & $<$ 0,001 kN \\
    Convergência & Momento máximo, 401 contra 3.201 pontos de integração (kN·m) & refino da discretização & 477,9 & 477,9 & $<$ 0,01~\% \\
    Numérico & Momento no meio do vão, carga média da combinação governante, sem beiral (kN·m) & OpenSeesPy, elementos de viga & 480,9 & 480,8 & 0,03~\% \\
\bottomrule
\end{tabular}}{Elaborado pelo autor (2026)}{referencias}{}{}


O equilíbrio é conferido em todas as combinações de estado-limite último: a soma
das reações reproduz a soma das cargas aplicadas à viga e ao balanço do beiral. A
convergência da integração ao longo do vão é conferida pelo refino de 401 para
3.201 pontos, que altera o momento máximo em menos de $10^{-5}$.

Dois modelos de elementos finitos, construídos no OpenSeesPy de forma
independente do integrador, conferem hipóteses do cálculo. O da viga representa
meia viga, por simetria, com 40 elementos \textit{elasticBeamColumn} ao longo do
eixo baricêntrico inclinado, cada um com a área e a inércia da seção no seu ponto
médio, rótula no apoio e engaste de simetria na cumeeira, e carga uniforme por
elemento; reproduz o momento do integrador com diferença de 0,03\%
(\autoref{tab:referencias}). O do diafragma da cobertura, em planta, tem nós nos
cruzamentos de terças e vigas, todas as barras rotuladas, cruzes só à tração e os
painéis dos frontões e verticais como molas; ele demonstra que sem as cruzes a
grelha é um mecanismo e fornece a rigidez transversal do conjunto
(\autoref{sec:ordem-analise}). Nenhum dos dois substitui uma verificação por
programa comercial de terceiros, que permanece recomendada antes da execução da
obra.

\section{ESTILO DO CÓDIGO}\label{sec:ap-estilo}

A \autoref{lst:ap-geometria} reproduz um trecho representativo, escolhido pela
brevidade e por ilustrar as convenções adotadas: documentação em português,
registro da expressão normativa aplicada e emprego de operações vetoriais.

\filbreak
\begin{lstlisting}[
    caption={Cálculo da altura da seção ao longo do vão},
    label={lst:ap-geometria}]
def h(self, x: float | NDArray) -> float | NDArray:
    """
    Altura da seção em x (m).

    A viga de duplo-bisel tem h crescente do apoio até a cumeeira.
    Para 0 <= x <= L/2: h(x) = h0 + x*tan(alpha)
    Por simetria: h(x) = h(L-x) para L/2 < x <= L
    """
    x = np.asarray(x, dtype=float)
    x_sim = np.where(x <= self.L / 2, x, self.L - x)
    return self.h0 + x_sim * np.tan(self.alpha_rad)

def A(self, x: float | NDArray) -> float | NDArray:
    """Área da seção transversal em x (m2)."""
    return self.b * self.h(x)

def I(self, x: float | NDArray) -> float | NDArray:
    """Momento de inércia em torno do eixo y em x (m4)."""
    return self.b * self.h(x) ** 3 / 12
\end{lstlisting}


\section{DISPONIBILIDADE E REPRODUÇÃO}\label{sec:ap-disponibilidade}

O código-fonte, a suíte de testes e os arquivos de saída deste trabalho,
incluindo o resultado integral da varredura paramétrica da
\autoref{sec:varredura-global}, estão no repositório indicado a seguir.

\vspace{\onelineskip}
\noindent
\url{https://github.com/gabstocki/mlc_projeto}
\vspace{\onelineskip}

A versão que produziu esta monografia é identificada pelo programa na própria
execução, e o registro abaixo é gerado junto com o memorial: \textit{commit}
\texttt{\versaoCommit}, de \versaoData, etiqueta \versaoTag, \versaoEstado.
Os dados de entrada são as constantes de projeto do módulo \texttt{main.py} e os
valores declarados da \autoref{tab:valores-declarados}. A reprodução
requer Python 3.12 e as dependências declaradas no arquivo
\texttt{pyproject.toml}, com os limites de versão nele registrados, e consiste
em três comandos, executados na raiz do repositório: a instalação do pacote, com
\texttt{pip install -e .[dev]}; a suíte de testes, com \texttt{pytest}; e o
\textit{pipeline}, com \texttt{python -m mlc\_projeto.main --sem-viz}, que
regenera o memorial, as tabelas, as figuras e os desenhos desta monografia.
